\chapter{Motivation und Zielsetzung}
\label{motivation-und-zielsetzung}

In diesem Kapitel werden die Beweggründe dieser Arbeit geschildert. Aus diesen Beweggründen 
werden dann die übergeordneten Ziele ausgearbeitet, welche als Leitlinie für diese Arbeit 
dienen. Diese Ziele sind Bestandteil der Forschungsfrage, aus der die 
Geschäftsanforderungen in Kapitel \ref{analyse-und-konzeption} abgeleitet werden, und 
bilden ebenfalls eine Grundlage für das Vorgehen in den nächsten Kapiteln.

\section{Motivation}

Wie in Kapitel \ref{einleitung} schon am Ende erwähnt wurde, gibt es eine Lücke zwischen dem Modell und der 
Ausführung bei Geschäftsprozessen. Der Soll-Zustand wird im Modell zwar beschrieben, aber der Ist-Zustand in der Produktion ist 
ein anderer. Ziel ist es, diese beiden Zustände durch diese Arbeit einander anzunähern, indem man durch eine Übersetzung 
des Prozessmodells in ein Ausführungsmodell eine Art Ablaufplan bekommt, der viel näher am Prozess liegt. Dies verspricht 
zum einen den Vorteil, dass alle Stakeholder einen klaren Überblick über das Prozessgeschehen haben. Bei der Produktion im Industriebereich
ist das ein starker Vorteil, da man als Produktionsleiter jeden Produktionsschritt auch anhand des Prozessdiagramms ablesen kann, anstatt hier ein
separates mentales Modell zu haben. Dadurch haben sowohl Management- als auch Produktionsebene dieselbe Modellbasis und können so den Produktionsprozess
leichter nachverfolgen und auch Optimierungen leichter durchführen. Diese Erleichterung in der Optimierung ist auch der zweite Vorteil. Viele Prozessverwaltungsplattformen
bieten Lösungen zur Prozessoptimierung an. Dabei handelt es sich meist um Operationen, die auf dem Prozessdiagramm oder auf dessen Basis stattfinden. Dadurch, dass das Prozessmodell nun auch Basis 
für die Ausführung ist, sind die Optimierungen auch viel näher am echten Geschehen und müssen nicht durch Richtlinien zunächst theoretisch festgelegt werden. Dadurch entfällt 
auch die Zwischenstufe von der Richtlinie zur tatsächlichen Umsetzung der Optimierung. Das spart Zeit und auch Ressourcen ein.
Auch wird dadurch das Produktionsgeschehen für alle transparenter, denn die Wissensbasis für das Prozessgeschehen ist nicht mehr eine Sammlung aus dem Diagramm abgeleiteter Dokumente, 
sondern das Prozessdiagramm selbst. Das Live-Streaming der Ausführung trägt zu diesem Aspekt am stärksten bei, da nun der laufende Produktionsprozess auf diesem Diagramm abgebildet wird.

\section{Zielsetzung}

In der heutigen Zeit steht die Digitalisierung von Produktionsprozessen mit Industrie 4.0 im Vordergrund. Jedoch gibt es hierbei Herausforderungen in der Interoperabilität\footnote{Die Fähigkeit, dass zwei oder mehrere Systeme Informationen austauschen und diese auch gegenseitig nutzen können. \cite{interoperability_iso}} durch fragmentierte Daten und Systeme \cite{huebner_semantic_2026}, Prozesseinhaltung durch Abweichungen des modellierten Prozesses \cite{van_der_aalst_replaying_2012} sowie Transparenz des Produktionsprozesses für Stakeholder \cite{jawad_adoption_2022}.

Aus dieser Ausgangsposition können folgende Ziele ausgearbeitet werden:

\begin{enumerate}
    \item Fabriken können heterogene Systeme einheitlich anbinden und ansteuern 
    \item Produktionen laufen prozesseinhaltend
    \item Der Produktionsprozess wird für alle Stakeholder transparent aufbereitet
\end{enumerate}

\subsection{Fabriken können heterogene Systeme einheitlich anbinden und ansteuern}

In den meisten Fällen hat man in einer Produktionsstätte ein Produktionssystem bestehend aus unterschiedlichen Maschinen. 
Diese haben dann eigene Schnittstellen, welche auch proprietär sein können, was zu Datensilos führen kann. Diese Arbeit soll durch Nutzung von 
offenen Kommunikations- und Schnittstellenstandards die Anbindung von heterogenen Systemen sowie die Ansteuerung erleichtern. Dies ist durch Schaffung einer einheitlichen Plattform möglich, 
welche die Übersetzung auf die unterschiedlichen Systeme übernimmt.

Um dieses Ziel zu erreichen, werden in dieser Arbeit folgende Arbeitsziele verfolgt:

\begin{itemize}
  \item \textbf{AZ1:} Auswahl eines offenen Kommunikationsprotokolls für die Ansteuerung von 
  IoT-Geräten und Definition einer einheitlichen Schnittstelle, über die jeder Prozessschritt als Endpunkt aufgerufen werden kann.

  \item \textbf{AZ2:} Automatische Generierung eines Firmware-Grundgerüsts für 
  Mikrocontroller aus dem Prozessmodell. Kommunikation, Netzwerkanbindung und Registrierung 
  der Endpunkte sind darin bereits enthalten, sodass Entwickler nur noch die 
  anlagenspezifische Logik implementieren.
\end{itemize}

\subsection{Produktionen laufen prozesseinhaltend}

In der Praxis laufen Prozesse nicht immer gemäß der Modellierung. Hierbei können die Gründe 
unterschiedlichster Art sein:

\begin{itemize}
  \item Menschliches Versagen
  \item Verschleiß
  \item Zustände der Umgebung
\end{itemize}

Durch das Nutzen des Prozessmodells als Basis für die automatische Prozessausführung wird 
sichergestellt, dass Produktionen prozessnäher funktionieren.

Dazu werden folgende Arbeitsziele umgesetzt:

\begin{itemize}
  \item \textbf{AZ3:} Formale Definition eines Prozessmodells auf Basis der EPK sowie eines 
  Ausführungsmodells, das die Prozessschritte und Verknüpfungsoperatoren in eine ausführbare 
  Form überführt.
  
  \item \textbf{AZ4:} Entwicklung eines Compilers, der ein EPK-Modell validiert und in das 
  Ausführungsmodell übersetzt.
  
  \item \textbf{AZ5:} Entwicklung einer Execution-Engine, die das Ausführungsmodell Schritt 
  für Schritt auf den angebundenen Geräten ausführt und Verzweigungen anhand der 
  Rückgabewerte gemäß dem Modell auflöst.
\end{itemize}

\subsection{Der Produktionsprozess wird für alle Stakeholder transparent aufbereitet}

Die Produktion findet in der Fabrikhalle statt. Hierbei können nicht alle Stakeholder einen 
Einblick haben. Dies wird derzeit beispielsweise durch Daten aus \ac{ERP}-Systemen ermöglicht. Hier 
wäre es also wichtig, näher am Prozessgeschehen zu sein, idealerweise in Echtzeit.

Dafür werden folgende Arbeitsziele verfolgt:

\begin{itemize}
  \item \textbf{AZ6:} Live-Darstellung der laufenden Ausführung direkt im Prozessmodell der 
  Prozessverwaltungsplattform.
  \item \textbf{AZ7:} Protokollierung jedes Ausführungsschritts in einem Ereignisprotokoll, 
  das für Process-Mining-Analysen exportiert werden kann.
\end{itemize}

Abschließend werden alle Komponenten in die Prozessverwaltungsplattform ARIS integriert (\textbf{AZ8}) und 
anhand eines Versuchsaufbaus validiert (\textbf{AZ9}).

\section{Vorgehen und Aufbau der Arbeit}

Nachdem in diesem Kapitel die Motivation sowie die übergeordneten Ziele und Arbeitsziele 
dargelegt wurden, legt Kapitel \ref{grundlagen-und-theoreitscher-rahmen} die theoretischen 
Grundlagen zur Lösung des Problems. In Kapitel \ref{related-work} werden verwandte Arbeiten betrachtet, 
welche ein ähnliches Problem lösen. Daraufhin wird ermittelt, welche Anteile dieser Arbeiten 
verwendet und welche neuen Beiträge geleistet werden. Als Nächstes werden in Kapitel 
\ref{analyse-und-konzeption} die Anforderungen an die Lösung der vorliegenden Arbeit definiert, 
zur Überprüfung in einem geschäftlichen Kontext. Basierend auf den Grundlagen von Kapitel 
\ref{grundlagen-und-theoreitscher-rahmen} sowie den Anforderungen aus Kapitel 
\ref{analyse-und-konzeption} wird in Kapitel \ref{systemdesign-und-spezifikation} das System 
zur Kompilierung und Ausführung von Prozessen formal entworfen. Darauf aufbauend findet in 
Kapitel \ref{technische-umsetzung} eine technische Umsetzung statt, welche das formal 
definierte System in der Praxis umsetzt. Anschließend erfolgt in Kapitel \ref{ergebnisse} 
eine Prüfung gegen die Anforderungen aus Kapitel \ref{analyse-und-konzeption} sowie eine 
Validierung durch die Anwendung des Prototyps der technischen Umsetzung in einem 
Beispielszenario. Folglich findet eine Diskussion bezüglich der Zielsetzung und der 
Umsetzung statt. Anschließend wird in Kapitel \ref{fazit} ein Fazit gezogen, gefolgt 
von einem zukünftigen Ausblick des Resultats der Arbeit sowie einer Bewertung der 
Praxistauglichkeit dieser Arbeit.
